\documentclass[a5paper,10pt]{article}

% https://github.com/InDevRus/filippov-solutions

% Нумерация страниц
\pagenumbering{gobble}

% Настройка полей
\usepackage{geometry}
\geometry{left=1cm}
\geometry{right=1cm}
\geometry{top=1cm}
\geometry{bottom=1.5cm}

% Вставка таблицы
\usepackage{array}
\usepackage{tabularx}

% Инструменты выравнивания формул
\usepackage{amsmath}
\usepackage{mathtools}

% Диакритика в математике
\usepackage{amsfonts}

% Вычисления размеров на ходу
\usepackage{calc}

% Удобные записи производных
\usepackage[italic=true]{derivative}

% Настройки сносок
\usepackage{enotez}

% Длинные знаки векторов
\usepackage{anyfontsize}
\usepackage[b]{esvect}

% Вставка картинок
\usepackage{graphicx}

% Вставка красной строки
\usepackage{indentfirst}
\setlength{\parindent}{1em}

% Условия в командах
\usepackage{xifthen}

% Луа-скрипты
\usepackage{luacode}

% Настройка команд отдельно в математическом и текстовом режимах
\usepackage{mathcommand}

% Для повторения LaTeX кода
\usepackage{multido}

% Конфигурация отступов формул
\delimitershortfall-1sp
\usepackage{mleftright}
\mleftright

% Растягивание объектов
\usepackage{relsize}

% Обтекание текстом
\usepackage{wrapfig}
\usepackage{subcaption}
\usepackage{floatflt}

% Поддержка ввода кириллицы
\usepackage[english, russian]{babel}

% Настройка корневой позиции для компилятора
\usepackage{import}
\providecommand*{\ProjectRootPath}{..}

\import{\ProjectRootPath/preambles/macros/}{theorems}

\import{\ProjectRootPath/preambles/fonts}{font_setup}

% Знак градуса
\usepackage{siunitx}

\import{\ProjectRootPath/preambles/macros/}{operators}
\import{\ProjectRootPath/preambles/macros/}{redefinitions}

\import{\ProjectRootPath/preambles/fonts}{letters_setup}
\import{\ProjectRootPath/preambles/fonts/adjustments}{custom_superscripts}

\import{\ProjectRootPath/preambles/custom}{formatting}
\import{\ProjectRootPath/preambles/custom}{frames}
\import{\ProjectRootPath/preambles/custom}{list_setup}

% TODO: перевести команды на современный стиль объявления

\begin{document}

$x\`\br{t} - a\br{x} x\br{t} = 0$.
Обозначим за $A\br{t}$ интеграл 
$A\br{t} = \tintlim{0}{t} a\br{s} \di s$, так что $\der{A}{t} = a\br{t}$.
$x\`\br{t} \pow{e}{- A\br{t}} - \pow{e}{- A\br{t}} a\br{x} x\br{t} = 0$.
$\der{}{t} \br{ x \pow{e}{- A\br{t}}} = 0$.
$x\br{t} = C \pow{e}{A\br{t}}$ -- общее решение.
Условие $\tuplim{\mathsmaller{t \to +\infty}} A\br{t} < +\infty$ равносильно ограниченности 
сверху этой функции, то есть тому, что существует такое $\alpha$, 
что $A\br{t} \le \alpha$ при всех $t \ge \sub{t}{0}$.

Далее за $\phi\br{t} = 0$ обозначено нулевое решение, $\sub{t}{0} = 0$.

\begin{necessity}
Пусть $A\br{t}$ сверху неограничена при $t \ge 0$. Покажем, что решение $\phi\br{t}$ неустойчиво.

Пусть $\epsilon = 1$, $\delta > 0$ -- произвольное и 
$\vbr{\x\br{0} - \phi\br{0}} = \vbr{C} < \delta$, $\vbr{C} \ne 0$. 
$A\br{t}$ как функция неограничена, значит найдётся такое $\sub{t}{1} > 0$,
что $A\br{\sub{t}{1}} \ge \ln\br{\dfrac {\epsilon} {\vbr{C}}}$.
В таком случае 
$\vbr{\x\br{\sub{t}{1}} - \phi\br{\sub{t}{1}}} 
    = \vbr{C} \, \pow{e}{A\br{\sub{t}{1}}} \ge \epsilon$.
Последнее означает, что решение неустойчиво.

\end{necessity}

\begin{sufficiency}

Пусть теперь $A\br{t}$ ограничена сверху числом $\alpha$. Покажем устойчивость решения.

Пусть $\epsilon > 0$ и $\delta\br{\epsilon} = \epsilon\pow{e}{-\alpha}$.
Потребуем, чтобы выполнялось $\vbr{\x\br{0} - \phi\br{0}} = \vbr{C} < \delta$. 
Из этого следует, что для всех $t \ge 0$
$$\vbr{\x\br{t} - \phi\br{t}} 
    = \vbr{C} \, \pow{e}{A\br{t}} 
    < \delta \pow{e}{\alpha} 
    = \epsilon.
$$

Таким образом, устойчивость решения установлена. Теорема полностью доказана.

\end{sufficiency}

\end{document}
